\documentclass{../../kalkulus-ppt}

\author[Tetew]{Teosofi Hidayah Agung}
\date{\today}
\title[Kalkulus 1 - Sup A]{Bilangan Kompleks}
\institute[Matematika ITS]{Departemen Matematika\\ Institut Teknologi Sepuluh Nopember}
\titlegraphic{{\includegraphics[scale=0.3]{ITS.png}$\quad$\includegraphics[scale=0.02]{M.png}}}

\newcommand{\dom}{\mathcal{D}}
\newcommand{\rng}{\mathcal{R}}

\renewcommand{\arraystretch}{1.5}

\begin{document}
\begin{frame}
    \titlepage
\end{frame}

\AtBeginSection{
    {
        \begin{frame}{Daftar isi}
            \tableofcontents[currentsection]
        \end{frame}}
}

\section{Bilangan Kompleks}
\begin{frame}
    \frametitle{\insertsection}
    \begin{definisi}[Bilangan Kompleks]
        Bilangan kompleks $z$ ditulis dalam bentuk
        \[z = a + bi\]
        di mana $a, b \in \mathbb{R}$ dan $i = \sqrt{-1}$ (satuan imajiner, sehingga $i^2 = -1$).\\
        $a$ disebut \textbf{bagian real} dari $z$, dinotasikan $\mathrm{Re}(z)$.\\
        $b$ disebut \textbf{bagian imajiner} dari $z$, dinotasikan $\mathrm{Im}(z)$.
    \end{definisi}
\end{frame}

\section{Teorema De Moivre}
\begin{frame}
    \frametitle{\insertsection}
    \begin{teorema}[Teorema De Moivre]
        Jika $z = r(\cos \theta + i \sin \theta)$ adalah bilangan kompleks dalam bentuk polar, maka untuk setiap bilangan bulat $n$, berlaku:
        \[z^n = r^n (\cos(n\theta) + i \sin(n\theta))\]
    \end{teorema}
\end{frame}

\end{document}
